\section{Математические модели механики жидкости и газа}
\label{sec:models-fluid-gas}

Механика жидкости и газа рассматривает движение и равновесие текучих
сплошных сред. Математическая модель строится из
\textbf{законов баланса массы, импульса и энергии}, определяющих
соотношений, уравнения состояния и начально-граничных условий.

Основные поля:
\[
\rho(\mathbf x,t),\qquad
\mathbf v(\mathbf x,t),\qquad
p(\mathbf x,t),\qquad
T(\mathbf x,t).
\]

\subsection{Гипотеза сплошной среды и кинематика}

В континуальном приближении молекулярная структура заменяется непрерывными
полями. Представительный физический объём должен быть мал по сравнению с
масштабом течения, но содержать большое число молекул. Критерием применимости
континуальной модели служит число Кнудсена
\[
Kn=\frac{\lambda}{L},
\]
где $\lambda$ --- длина свободного пробега, $L$ --- характерный размер. При
$Kn\ll1$ модель сплошной среды обычно применима.

Траектория жидкой частицы удовлетворяет
\[
\frac{d\mathbf x}{dt}=\mathbf v(\mathbf x,t).
\]
Для произвольного поля $\phi$ производная вдоль движения равна
\[
\boxed{
\frac{D\phi}{Dt}
=\frac{\partial\phi}{\partial t}
+(\mathbf v\cdot\nabla)\phi
}
\]
и называется \textbf{материальной производной}.

Градиент скорости раскладывается как
\[
\nabla\mathbf v=\mathbf D+\mathbf W,
\]
где
\[
\mathbf D=\frac12\left(\nabla\mathbf v+(\nabla\mathbf v)^T\right)
\]
--- тензор скоростей деформации, а антисимметричная часть $\mathbf W$
характеризует локальное вращение. Вихрь
\[
\boldsymbol\omega=\nabla\times\mathbf v.
\]

\subsection{Закон сохранения массы}

Для неподвижного контрольного объёма $V$
\[
\frac{d}{dt}\int_V\rho\,dV
+\int_{\partial V}\rho\mathbf v\cdot\mathbf n\,dS=0.
\]
С помощью формулы Гаусса получаем локальное уравнение неразрывности
\[
\boxed{
\frac{\partial\rho}{\partial t}
+\nabla\cdot(\rho\mathbf v)=0
}
\]
или
\[
\frac{D\rho}{Dt}+\rho\nabla\cdot\mathbf v=0.
\]
Для жидкости постоянной плотности
\[
\boxed{\nabla\cdot\mathbf v=0}.
\]

\subsection{Закон сохранения импульса}

Локальный баланс количества движения имеет вид
\[
\boxed{
\rho\frac{D\mathbf v}{Dt}
=\nabla\cdot\boldsymbol\sigma+\rho\mathbf f,
}
\]
где $\boldsymbol\sigma$ --- тензор напряжений, $\mathbf f$ --- массовая
сила на единицу массы. Для жидкости удобно выделять давление:
\[
\boldsymbol\sigma=-p\mathbf I+\boldsymbol\tau,
\]
поэтому
\[
\rho\frac{D\mathbf v}{Dt}
=-\nabla p+\nabla\cdot\boldsymbol\tau+\rho\mathbf f.
\]

\subsection{Закон сохранения энергии}

Пусть $e$ --- удельная внутренняя энергия, а
\[
E=e+\frac{|\mathbf v|^2}{2}
\]
--- полная удельная энергия. Уравнение внутренней энергии можно записать как
\[
\boxed{
\rho\frac{De}{Dt}
=-p\,\nabla\cdot\mathbf v
+\boldsymbol\tau:\mathbf D
-\nabla\cdot\mathbf q
+\rho r,
}
\]
где $\mathbf q$ --- тепловой поток, $r$ --- объёмный источник тепла на
единицу массы. Эквивалентная консервативная форма полной энергии:
\[
\frac{\partial(\rho E)}{\partial t}
+\nabla\cdot\bigl[(\rho E+p)\mathbf v\bigr]
=\nabla\cdot(\boldsymbol\tau\mathbf v)
-\nabla\cdot\mathbf q
+\rho\mathbf f\cdot\mathbf v+\rho r.
\]

\subsection{Определяющие соотношения и замыкание}

Законы сохранения универсальны, но для получения замкнутой системы нужны
свойства конкретной среды. Для ньютоновской жидкости
\[
\boxed{
\boldsymbol\tau
=2\mu\mathbf D
+\lambda_v(\nabla\cdot\mathbf v)\mathbf I,
}
\]
где $\mu$ и $\lambda_v$ --- коэффициенты вязкости. Для несжимаемой
ньютоновской жидкости
\[
\boldsymbol\tau=2\mu\mathbf D.
\]
Закон Фурье:
\[
\boxed{\mathbf q=-k\nabla T}.
\]
Для идеального газа
\[
\boxed{p=\rho RT},
\qquad
e=c_vT.
\]
Именно выбор определяющих соотношений отличает ньютоновские и
неньютоновские жидкости, идеальный газ и более сложные среды.

\subsection{Уравнения Эйлера и Навье--Стокса}

Если вязкостью и теплопроводностью пренебречь, получаются уравнения Эйлера.
Для сжимаемого потока:
\[
\frac{\partial\rho}{\partial t}+\nabla\cdot(\rho\mathbf v)=0,
\]
\[
\frac{\partial(\rho\mathbf v)}{\partial t}
+\nabla\cdot(\rho\mathbf v\otimes\mathbf v+p\mathbf I)
=\rho\mathbf f.
\]
Уравнение энергии дополняет систему.

Для несжимаемой ньютоновской жидкости постоянной плотности система
Навье--Стокса принимает вид
\[
\boxed{
\begin{aligned}
\nabla\cdot\mathbf v&=0,\\
\frac{\partial\mathbf v}{\partial t}
+(\mathbf v\cdot\nabla)\mathbf v
&=-\frac1\rho\nabla p+\nu\Delta\mathbf v+\mathbf f,
\end{aligned}
}
\]
где $\nu=\mu/\rho$ --- кинематическая вязкость. Нелинейность связана с
конвективным членом $(\mathbf v\cdot\nabla)\mathbf v$.

Для стационарного невязкого течения жидкости постоянной плотности под
действием консервативной массовой силы $\mathbf f=-\nabla\Pi$ справедлива
\textbf{теорема Бернулли}:
\[
\boxed{
\frac{p}{\rho}+\frac{|\mathbf v|^2}{2}+\Pi=\text{const}
}
\]
вдоль линии тока. В поле тяжести при вертикальной координате $z$, направленной
вверх, $\Pi=gz$, поэтому
\[
\frac{p}{\rho}+\frac{|\mathbf v|^2}{2}+gz=\text{const}.
\]
При безвихревом течении константа одинакова во всей связной области.

\subsection{Начальные и граничные условия}

Для нестационарной задачи задают начальные поля, например
\[
\mathbf v(\mathbf x,0)=\mathbf v_0(\mathbf x),
\qquad
\rho(\mathbf x,0)=\rho_0(\mathbf x).
\]
На неподвижной твёрдой стенке для вязкой жидкости используется условие
прилипания
\[
\mathbf v=0.
\]
Для идеальной жидкости обычно достаточно непротекания
\[
\mathbf v\cdot\mathbf n=0.
\]
На входе и выходе задают совместимый набор скоростей, расходов, давления,
температуры и других характеристик. Для сжимаемых течений число допустимых
граничных условий связано с направлениями характеристик.

\subsection{Безразмерные критерии}

После введения характерных масштабов $L$, $U$, $\rho_0$ у несжимаемых
уравнений Навье--Стокса появляется число Рейнольдса
\[
\boxed{Re=\frac{\rho_0UL}{\mu}=\frac{UL}{\nu}},
\]
характеризующее отношение инерционных и вязких эффектов. Для сжимаемых
потоков важно число Маха
\[
Ma=\frac{U}{a},
\]
для течений со свободной поверхностью --- число Фруда, а при совместном
переносе импульса и тепла --- числа Прандтля и Пекле. Безразмерные числа
определяют режимы течения и позволяют сравнивать физически разные системы.

\subsection{Турбулентность и численное моделирование}

При больших числах Рейнольдса движение может стать турбулентным. Основные
уровни моделирования:
\begin{itemize}
    \item DNS --- непосредственное разрешение всех существенных масштабов;
    \item LES --- разрешение крупных вихрей с моделью подсеточных масштабов;
    \item RANS --- осреднённые уравнения с моделью турбулентных напряжений.
\end{itemize}
Это пример \textbf{проблемы замыкания}: после осреднения возникают новые
неизвестные корреляции, для которых нужны дополнительные модели.

Численно применяются методы конечных разностей, конечных объёмов и конечных
элементов. Для законов сохранения особенно важна консервативность
дискретизации; для задач с сильной конвекцией --- устойчивость и корректное
описание разрывов.


\subsection{Интегральная форма законов сохранения}

Локальные уравнения механики жидкости и газа получаются из балансов для
контрольного объёма $V$ с границей $\partial V$. Например, закон сохранения массы
имеет вид
\[
    \boxed{
    \frac{d}{dt}\int_V\rho\,dV
    +\int_{\partial V}\rho\mathbf v\cdot\mathbf n\,dS=0.
    }
\]
Применение теоремы Гаусса даёт локальное уравнение неразрывности. Интегральная
форма особенно важна для методов конечных объёмов и для решений с разрывами,
где классические производные могут не существовать.

Для импульса аналогично
\[
    \frac{d}{dt}\int_V\rho\mathbf v\,dV
    +\int_{\partial V}\rho\mathbf v(\mathbf v\cdot\mathbf n)\,dS
    =\int_{\partial V}\boldsymbol\sigma\mathbf n\,dS
     +\int_V\rho\mathbf f\,dV.
\]

\subsection{Консервативная форма уравнений Эйлера для газа}

Для идеального сжимаемого газа без вязкости и теплопроводности удобно ввести
полную удельную энергию
\[
    E=e+\frac{|\mathbf v|^2}{2}.
\]
Тогда система Эйлера записывается в консервативной форме
\[
\boxed{
\begin{aligned}
\partial_t\rho+\nabla\cdot(\rho\mathbf v)&=0,\\
\partial_t(\rho\mathbf v)
+\nabla\cdot(\rho\mathbf v\otimes\mathbf v+p\mathbf I)&=\rho\mathbf f,\\
\partial_t(\rho E)
+\nabla\cdot\bigl((\rho E+p)\mathbf v\bigr)&=\rho\mathbf f\cdot\mathbf v.
\end{aligned}}
\]
Система замыкается уравнением состояния, например
\[
    p=(\gamma-1)\rho e.
\]
Консервативная запись принципиальна при ударных волнах и других разрывах.

\subsection{Условия Ранкина--Гюгонио}

Если через поверхность разрыва, движущуюся с нормальной скоростью $s$, проходит
консервативная величина
\[
    \partial_t U+\nabla\cdot F(U)=0,
\]
то слабое решение должно удовлетворять условию скачка
\[
    \boxed{s[U]=[F_n(U)]},
\]
где $[\cdot]$ обозначает разность значений по разные стороны поверхности, а
$F_n=F\cdot n$. Для газовой динамики эти соотношения выражают сохранение массы,
импульса и энергии на ударной волне. Одних условий сохранения недостаточно для
выбора физического разрыва; дополнительно используется условие энтропии.

\subsection{Безразмерная форма Навье--Стокса}

При масштабах $L$ и $U$ положим
\[
    \mathbf x=L\widetilde{\mathbf x},
    \qquad
    t=\frac{L}{U}\widetilde t,
    \qquad
    \mathbf v=U\widetilde{\mathbf v},
    \qquad
    p=\rho U^2\widetilde p.
\]
Для несжимаемой жидкости получаем
\[
    \boxed{
    \partial_{\widetilde t}\widetilde{\mathbf v}
    +(\widetilde{\mathbf v}\cdot\widetilde\nabla)\widetilde{\mathbf v}
    =-\widetilde\nabla\widetilde p
    +\frac1{Re}\widetilde\Delta\widetilde{\mathbf v}
    +\widetilde{\mathbf f}.
    }
\]
Поэтому $Re$ буквально является коэффициентом, определяющим относительный вес
вязкого члена по сравнению с инерционными членами.

\subsection{Вихрь и перенос вихря}

Вихрь скорости определяется как
\[
    \boldsymbol\omega=\nabla\times\mathbf v.
\]
Для несжимаемой жидкости с постоянной вязкостью при консервативных массовых силах
из Навье--Стокса следует
\[
    \boxed{
    \frac{D\boldsymbol\omega}{Dt}
    =(\boldsymbol\omega\cdot\nabla)\mathbf v
    +\nu\Delta\boldsymbol\omega.
    }
\]
Первое слагаемое --- растяжение вихревых линий; оно отсутствует в строго
двумерном течении. Второе описывает вязкую диффузию вихря. Это уравнение хорошо
показывает физическую роль нелинейности и вязкости.

\subsection{Осреднение Рейнольдса и проблема замыкания}

В RANS-подходе поле скорости раскладывают как
\[
    \mathbf v=\overline{\mathbf v}+\mathbf v',
    \qquad \overline{\mathbf v'}=0.
\]
После осреднения нелинейного конвективного члена появляется тензор
\textbf{рейнольдсовых напряжений}
\[
    -\rho\,\overline{v_i'v_j'}.
\]
Он содержит новые неизвестные корреляции и не выражается непосредственно через
среднее поле. Поэтому нужны дополнительные модели турбулентности --- это и есть
классическая проблема замыкания. В LES аналогичная проблема переносится на
подсеточные масштабы, а DNS избегает модели турбулентности ценой очень высокой
вычислительной стоимости.

\subsection{Бернулли для баротропной сжимаемой среды}

Для стационарного невязкого баротропного течения, где $\rho=\rho(p)$, вдоль
линии тока выполняется более общая форма
\[
    \boxed{
    \frac{|\mathbf v|^2}{2}
    +\int^p\frac{dp'}{\rho(p')}
    +\Pi=\text{const}.
    }
\]
При постоянной плотности интеграл превращается в $p/\rho$, и получается формула,
приведённая выше. Это подчёркивает, что простая запись $p/\rho+v^2/2+\Pi$ не
должна безоговорочно переноситься на произвольное сжимаемое течение.

\subsection{Что важно сказать на экзамене}

Ответ удобно строить как цепочку
\[
\boxed{
\text{сплошная среда}
\to
\text{масса, импульс, энергия}
\to
\text{определяющие соотношения}
\to
\text{граничные условия}
}.
\]
Обязательно записать уравнение неразрывности и систему Навье--Стокса для
несжимаемой жидкости, объяснить роль давления и вязкости. Полезно привести
теорему Бернулли с её условиями применимости и число Рейнольдса.

\newpage
